% !TEX root = ../main.tex
\chapter{O processador SAPHO}
\label{cap:processador-sapho}

Antes de criar processadores na IDE, vale entender o que exatamente está sendo criado. Este capítulo descreve a arquitetura do \emph{soft-processor} SAPHO do ponto de vista de quem o usa: o que os parâmetros significam, como o processador executa o seu programa e quais são as suas restrições. Não é preciso decorar nada aqui para usar a plataforma --- mas quem entende a máquina escreve programas melhores para ela.

\section{Visão geral da arquitetura}

O SAPHO é um processador de \textbf{acumulador} com \textbf{arquitetura Harvard}:

\begin{itemize}
  \item \textbf{Acumulador}: em vez de um banco de registradores, existe um único registrador central, o acumulador, que recebe o resultado de toda operação da ULA (unidade lógico-aritmética). O padrão dominante é: um operando vem da memória de dados, o outro é o próprio acumulador, e o resultado volta para o acumulador. Uma \textbf{pilha de dados} fornece o segundo operando quando a expressão exige.
  \item \textbf{Harvard}: as memórias de programa e de dados são fisicamente separadas, cada uma com sua largura e seu barramento. O programa não pode se sobrescrever, e dados e instruções são acessados em paralelo.
\end{itemize}

Três traços definem o estilo do processador:

\begin{enumerate}
  \item \textbf{Totalmente parametrizável} --- toda largura é um parâmetro Verilog: a palavra de dados (\code{NUBITS}), a mantissa (\code{NBMANT}) e o expoente (\code{NBEXPO}) do ponto flutuante, o operando das instruções, as profundidades das memórias e das pilhas.
  \item \textbf{Pague pelo que usar} --- cada instrução do conjunto é um parâmetro booleano; o montador do YANC liga apenas as instruções que o seu programa realmente usa, e o hardware das demais é eliminado na geração do circuito.
  \item \textbf{Ciclo curto} --- o processador opera com um \emph{pipeline} raso de três estágios (busca, decodificação, execução), sem \emph{forwarding}: a ULA é combinacional e o resultado se acumula a cada ciclo, de modo que nem saltos condicionais quebram o fluxo. O comportamento temporal é simples de prever, o que é uma virtude em instrumentação.
\end{enumerate}

\section{Os parâmetros fundamentais}
\label{sec:parametros-fundamentais}

Ao criar um processador (Capítulo~\ref{cap:criando-processadores}) você escolhe, direta ou indiretamente, os valores abaixo. Eles aparecem como diretivas no topo do arquivo \cmm{} e como \code{parameter} no Verilog gerado.

\begin{longtable}{@{}llp{0.60\linewidth}@{}}
\toprule
\textbf{Parâmetro} & \textbf{Diretiva \cmm{}} & \textbf{Significado} \\
\midrule
\endhead
\code{NUBITS} & \code{\#NUBITS} & Largura da palavra de dados, em bits. Todo inteiro do programa tem esse tamanho, em complemento de dois. \\
\code{NBMANT} & \code{\#NBMANT} & Largura da mantissa do ponto flutuante. \\
\code{NBEXPO} & \code{\#NBEXPO} & Largura do expoente do ponto flutuante. \\
--- & \code{\#NDSTAC} & Profundidade da pilha de dados. \\
--- & \code{\#SDEPTH} & Profundidade da pilha de sub-rotinas (chamadas de função aninhadas). \\
--- & \code{\#NUIOIN} & Número de portas de entrada. \\
--- & \code{\#NUIOOU} & Número de portas de saída. \\
--- & \code{\#NUGAIN} & Constante usada pela função \code{norm()} (normalização por divisão fixa). \\
--- & \code{\#FFTSIZ} & Tamanho da FFT ($2^{\texttt{FFTSIZ}}$ pontos), quando o programa usa os endereçamentos com reversão de bits. \\
\bottomrule
\end{longtable}

\begin{aviso}
\textbf{Restrição estrutural do ponto flutuante:} a palavra precisa comportar exatamente o formato, isto é, \texttt{NUBITS} = \texttt{NBMANT} + \texttt{NBEXPO} + 1 (o bit extra é o sinal). A AURORA valida essa igualdade no formulário de criação de processador e recusa combinações inválidas.
\end{aviso}

\section{O ponto flutuante do SAPHO}
\label{sec:float-sapho}

O SAPHO tem hardware de ponto flutuante \textbf{próprio, que não segue o padrão IEEE 754}. O formato é: 1 bit de sinal, \code{NBEXPO} bits de expoente em complemento de dois e \code{NBMANT} bits de mantissa com o bit líder explícito. Isso permite escolher precisões arbitrárias --- uma mantissa de 10 bits com expoente de 5, por exemplo --- e economizar hardware exatamente onde a aplicação permite.

Na prática, para o usuário:
\begin{itemize}
  \item Dentro do programa \cmm{}, \code{float} \enquote{simplesmente funciona}: literais, aritmética e as funções da biblioteca operam nesse formato interno.
  \item A conversão entre o formato IEEE do seu computador e o formato do SAPHO é feita em software pela toolchain (na carga das memórias e na leitura dos resultados), de forma transparente.
  \item A precisão que você obtém é exatamente a que você configurou: com \code{\#NBMANT 10}, espere cerca de 3 dígitos decimais significativos.
\end{itemize}

Números \textbf{complexos} (tipo \code{comp} da linguagem) não têm hardware dedicado: cada complexo é um par de \emph{floats} (parte real e imaginária) manipulado componente a componente pelo código gerado.

\section{Como o programa vira circuito}

O processador gerado é composto por poucos módulos Verilog, todos criados a partir de moldes parametrizáveis da pasta de HDL do YANC:

\begin{longtable}{@{}lp{0.68\linewidth}@{}}
\toprule
\textbf{Módulo} & \textbf{Função} \\
\midrule
\endhead
\code{processor} & O topo: instancia o núcleo e as duas memórias. \\
\code{mem\_instr} & Memória de programa (somente leitura), carregada com a imagem \file{<proc>_inst.mif}. \\
\code{mem\_data} & Memória de dados, carregada com a imagem \file{<proc>_data.mif}; uma leitura e uma escrita por ciclo. \\
\code{core} & O caminho de dados: acumulador, pilhas, busca de instruções, contador de programa. \\
\code{instr\_dec} & Decodificador de instruções. \\
\code{ula} & A unidade lógico-aritmética: um submódulo por operação, e só os usados sobrevivem. \\
\code{addr\_dec} & Decodificador de endereços para múltiplas portas de E/S. \\
\code{myFIFO} & FIFO genérica, disponível para acoplar o processador ao mundo externo sem depender de IP de fabricante. \\
\bottomrule
\end{longtable}

O conjunto de instruções usa opcode de 7 bits, com famílias para carga e armazenamento (direto, indireto e com reversão de bits para FFT), pilha, entrada e saída, controle de fluxo (salto incondicional, salto se zero, chamada e retorno de sub-rotina) e cálculo (aritmética inteira e em ponto flutuante, lógica, comparações, deslocamentos e conversões). Você não escreve essas instruções à mão --- o compilador \cmm{} as gera ---, mas as verá nas formas de onda: a AURORA exibe, em \emph{lockstep} com o clock, a instrução assembly e a linha \cmm{} correspondente a cada ciclo (Capítulo~\ref{cap:ondas}).

\section{Entrada, saída e o mundo externo}

O processador conversa com o exterior por portas numeradas de entrada e de saída, criadas conforme \code{\#NUIOIN} e \code{\#NUIOOU}. No programa, elas são acessadas pelas funções \code{in()}, \code{out()}, \code{fin()} e \code{fout()} (Capítulo~\ref{cap:cmm-io-stdlib}). No hardware, cada porta vira sinais de dados e de controle no topo do módulo \code{processor} --- é por eles que o seu projeto de FPGA conecta o processador a ADCs, DACs, FIFOs e outros blocos.

Dois recursos opcionais completam a interface: \textbf{interrupção} (o processador desvia para um endereço configurado quando o pino de interrupção pulsa) e o marcador \code{\#TOAQUI}, que gera um pino \code{cheguei} pulsado quando o programa passa por um ponto marcado --- útil para sincronizar hardware externo com fases do programa.

\section{Um processador por tarefa}

Como cada processador é gerado sob medida, o padrão de projeto no SAPHO é diferente do mundo dos processadores fixos: em vez de um processador grande que faz tudo, criam-se \textbf{vários processadores pequenos}, um por tarefa, conectados entre si por um \emph{top-level} Verilog (por exemplo, um processador que detecta cruzamentos de zero alimentando outro que calcula uma distância). A AURORA suporta esse fluxo diretamente: um projeto pode conter quantos processadores você quiser, cada um com sua pasta, seu \file{.cmm} e seus parâmetros, todos integrados por um arquivo \emph{top-level} e um \emph{testbench} comuns (Capítulo~\ref{cap:projetos}).
